\documentclass[spanish,openany]{report}
% \usepackage{darkmode} \enabledarkmode
\usepackage{wrapfig}
\input{~/git/Clase/template/preamble.tex}
\input{~/git/Clase/template/macros.tex}
\input{~/git/Clase/template/letterfonts.tex}
\graphicspath{{/home/kuma/git/Clase/Fp_Gs_Robotica/images/}}
\setimagefolder{/home/kuma/git/Clase/Fp_Gs_Robotica/images/}


\date{\today}
% \hbadness=99999
\usepackage{microtype}
% \addbibresource{./bibliografia.bib}
\usepackage{titling}
\usepackage{afterpage}


% Header and footer
% \fancypagestyle{plain}{
%     \fancyhf{} %Clear Everything.
%     \fancyfoot[C]{Página \thepage\ de \pageref{LastPage} }%Page Number
%     \renewcommand{\headrulewidth}{1pt} %0pt for no rule, 2pt thicker etc...
%     \renewcommand{\footrulewidth}{1pt}
%     \fancyfoot[R]{\small{Daniel Carbajales Soto}}
%     \fancyhead[R]{\small{Cinta Transportadora con 3 Velocidades}}
%     \fancyhead[L]{\small{\today}}
% }
% \pagestyle{plain}

\begin{document}

%----------------------------------------------------------------------------------------
%	                            PORTADA
%----------------------------------------------------------------------------------------
% \makeportada{}{Daniel Carbajales Soto}{ 1° CFGS Automatización y Robótica Industrial}{}

%----------------------------------------------------------------------------------------
%                                    HEADER 
%----------------------------------------------------------------------------------------
\header{SEHN - Sistemas Eléctricos Hidráulicos y Neumáticos }{Daniel Carbajales Soto}{1° CFGSARI}
%----------------------------------------------------------------------------------------
%                               TABLE OF CONTENTS
%----------------------------------------------------------------------------------------
% \maketoc
%----------------------------------------------------------------------------------------
%	                            TABLES
%----------------------------------------------------------------------------------------
% this is just a demo in case i need to make a table
%
%
%
% \begin{tabularx}{\linewidth}{|C{2cm}|Y|Y|C{2cm}|}
% \hline
% \makecell{Header 1} & \makecell{Header 2 \\ with linebreak} & \makecell{Header 3} & \makecell{Header 4} \\
% \hline
% Row 1 & This is a long cell text that automatically wraps in the Y column & Another long cell & Short \\
% \hline
% Row 2 & \multirow{2}{*}{Multirow text} & More text & Short \\
% \cline{1-1} \cline{3-4}
% Row 3 & & Extra text & Short \\
% \hline
% \end{tabularx}


%----------------------------------------------------------------------------------------
%	                            DOCUMENT
%----------------------------------------------------------------------------------------



\begin{enumerate}
	\item
		\begin{quote}
  		% ✍️ Write your response here
		\end{quote}
		
	\item
		\begin{quote}
  		% ✍️ Write your response here
		\end{quote}
		
	\item
		\begin{quote}
  		% ✍️ Write your response here
		\end{quote}
		
	\item
		\begin{quote}
  		% ✍️ Write your response here
		\end{quote}
		
	\item
		\begin{quote}
  		% ✍️ Write your response here
		\end{quote}
		
	\item\boldit{Dos corrientes alternas senoidales están desfasadas 20º. Sabiendo que la frecuencia es de 50 Hz. Calcular:}

		\begin{quote}
  		% ✍️ Write your response here
			\begin{tasks}(2)
				\task El período
				\task El tiempo de desfase
			\end{tasks}

			\begin{paracol}{2}
				\incfig[0.65]{desfase-tiempo}
			\switchcolumn
			\qs{}{ 
			a) El período es de: 
			\[
			\begin{split}
				t &= \frac{1}{50\si{\hertz}} \\
				&= \boxed{0'02\si{\second}}
			\end{split}
			\]
			}
			\end{paracol}
			\qs{}{
			b) El tiempo de desfase es de:
			\[\begin{split}
				t_{\text{desfase}} &= \frac{t}{\ang{360}}\times\ang{20} \\
				\frac{\frac{1}{50}}{\ang{360}}\times\ang{20} &= 0'0011\si{\second}
			\end{split}\]
			}
		\end{quote}

		
	\item
		\begin{quote}
  		% ✍️ Write your response here
		\end{quote}
		
\clearpage
	\item
		\begin{quote}
  		% ✍️ Write your response here
		\end{quote}
		
	\item
		\begin{quote}
  		% ✍️ Write your response here
		\end{quote}
		
	\item
		\begin{quote}
  		% ✍️ Write your response here
		\end{quote}
		
	\item
		\begin{quote}
  		% ✍️ Write your response here
		\end{quote}
		
	\item
		\begin{quote}
  		% ✍️ Write your response here
		\end{quote}
		
	\item\textbf{\textit{{Una bobina de resistencia $R = 200 , \Omega$ y coeficiente de autoinducción $L = 0,2 , \mathrm{H}$ se conecta en serie con un condensador de capacidad $C = 100 , \mu\mathrm{F}$ a una tensión alterna senoidal de $250 , \mathrm{V}$, $50 , \mathrm{Hz}$. Calcular:}}}
		\begin{quote}
  		% ✍️ Write your response here
		\begin{enumerate}[label=\alph*)]
				\item Reactancia total.
				\item Impedancia del circuito.
				\item Intensidad de corriente.
				\item Ángulo de desfase entre la tensión aplicada al circuito y la intensidad.
				\item Potencia activa, reactiva y aparente del circuito.
			\end{enumerate}
			\textbf{Solución:}
			\begin{enumerate}[label=\alph*)]
				\item $31 , \Omega$;
				\item $202,39 , \Omega$;
				\item $1,24 , \mathrm{A}$;
				\item $8^\circ 41'$;
				\item $307,3 , \mathrm{W}$, $47,6 , \mathrm{VAR}$ y $310 , \mathrm{VA}$.
			\end{enumerate}

		\end{quote}
		
	\item\boldit{En el circuito de la figura calcular la indicación de los aparatos de medida.}
		\begin{quote}
		% ✍️ Write your response here
		\begin{paracol}{2}
		
			\begin{data}
			\begin{center}
				\begin{tabular*}{\columnwidth}{@{\extracolsep{\fill}} c c}

					\hline
					\textbf{Magnitud} & \textbf{Valor}  \\ \hline\hline
					$U$ & $230$  \si{\volt}  \\ \hline
					$R$ & $80$  \si{\ohm}  \\ \hline
					$L$ & $0,01$  \si{\henry}  \\ \hline
					$C$ & $100$  \si{\micro\farad}  \\ \hline
					$f$ & $50$  \si{\hertz}  \\ \hline
				\end{tabular*}
			\end{center}
			\end{data}
			\switchcolumn
			\begin{solution}[Amperaje]
			\[
				\begin{split}
				I &= \frac{U}{R} \\
				I &= \frac{230}{80} \\
				I &= \boxed{2.7 \si{\ampere}} \\
				\end{split}
			\]
			\end{solution}
		\end{paracol}

			\begin{solution}[$U_{ab}$]
				\begin{align*}
				\[U = I Z \] 
				\[U = 2.7\cdot\sqrt{80^2 + (\frac{1}{2\pi\cdot 50 \cdot (100\cdot 10^{-6})})^2}\]
				\[U = \boxed{1.8 \si{\volt}} \]
				\end{align*}
			\end{solution}
			\begin{solution}[$U_{bc}$]
			\end{solution}


		\end{quote}
		
	\clearpage
	\item\boldit{En el circuito de la figura calcular la indicación de los aparatos de medida.}
		\begin{quote}
		% ✍️ Write your response here
		\end{quote}
	\item\boldit{Una bobina de resistencia 30$\Omega$  y coeficiente de autoinducción 0,4 H está conectada en serie con un condensador de 40 F a una tensión alterna senoidal de 220 V, 50 Hz.}



	\boldit{Calcular:}
		\begin{quote}
		% ✍️ Write your response here
			\begin{enumerate}[label=\alph*)]

				\item Triángulo de resistencias.


					\begin{paracol}{2}
						\[
							\begin{split}
								X_L &= 2\pi f 0,4(L) \\ 
								    &= 125,665\Omega
							\end{split}
						\]
						\switchcolumn
						\[
							\begin{split}
							X_c &= \frac{1}{2\pi \cdot 50 \cdot (40 \cdot 10^{-6})} \\
							    &= 79{,}775\Omega 
							\end{split}
						\]
						\switchcolumn
						\[
						\begin{split}
						X &= X_L - X_C \\
						  &= 125{,}664 \, \Omega - 79{,}577 \, \Omega \\
						  &= 46{,}087 \, \Omega
						\end{split}
						\]
						\vspace{1em}
						\[Zx = \sqrt{R^2 + X^2} \quad ;\quad Z_x = \sqrt{30^2 + 46,086^2} = 54,99\Omega \] 
						\switchcolumn
						\hfill\incfig[0.45]{Triangulo_resistencias}\hfill
					\end{paracol}

				\item Intensidad de corriente.

				\item Triángulo de tensiones.

				\item Triángulo de potencias.

			\end{enumerate}

		\end{quote}
	\boldit{Solución: }
	\begin{quote}
		\begin{enumerate}[label=\alph*)]
			\item Z = 55 \si{\ohm};
			\item 4 A;
			\item V$_x$ = 184 V y V$_R$ = 120 V.
			\item 480 W, 736 VAR y 880 VA.
			\end{enumerate}
	\end{quote}
	\item
		\begin{quote}
  		% ✍️ Write your response here
		\end{quote}
		
	\item
		\begin{quote}
  		% ✍️ Write your response here
		\end{quote}
		
	\item
		\begin{quote}
  		% ✍️ Write your response here
		\end{quote}
		
	\item
		\begin{quote}
  		% ✍️ Write your response here
		\end{quote}
		
	\item
		\begin{quote}
  		% ✍️ Write your response here
		\end{quote}
		
	\item
		\begin{quote}
  		% ✍️ Write your response here
		\end{quote}
		
	\item
		\begin{quote}
  		% ✍️ Write your response here
		\end{quote}
		
	\item
		\begin{quote}
  		% ✍️ Write your response here
		\end{quote}
		
	\item
		\begin{quote}
  		% ✍️ Write your response here
		\end{quote}
		
	\item
		\begin{quote}
  		% ✍️ Write your response here
		\end{quote}
		
	\item
		\begin{quote}
  		% ✍️ Write your response here
		\end{quote}
		
	\item
		\begin{quote}
  		% ✍️ Write your response here
		\end{quote}
		
	\item
		\begin{quote}
  		% ✍️ Write your response here
		\end{quote}
		
	\item
		\begin{quote}
  		% ✍️ Write your response here
		\end{quote}
		
	\item
		\begin{quote}
  		% ✍️ Write your response here
		\end{quote}
		
	\item
		\begin{quote}
  		% ✍️ Write your response here
		\end{quote}
		
	\item
		\begin{quote}
  		% ✍️ Write your response here
		\end{quote}
		
	\item
		\begin{quote}
  		% ✍️ Write your response here
		\end{quote}
		
	\item
		\begin{quote}
  		% ✍️ Write your response here
		\end{quote}
		
	\item
		\begin{quote}
  		% ✍️ Write your response here
		\end{quote}
		
	\item
		\begin{quote}
  		% ✍️ Write your response here
		\end{quote}
		
	\item
		\begin{quote}
  		% ✍️ Write your response here
		\end{quote}
		
	\item
		\begin{quote}
  		% ✍️ Write your response here
		\end{quote}
		
\end{enumerate}


%%%%%%%%%%%%%%%%%%%%%%%%%%%%%%%%%%%%%%%%%%%%%%%%%%%%%%%%%%
%%%%%%%%%%%%%%%%%%%% bibliography %%%%%%%%%%%%%%%%%%%%%%%%
% \nocite{*}                                 %%%%%%%%%%%%%
% \printbibliography                         %%%%%%%%%%%%%
%%%%%%%%%%%%%%%%%%%%%%%%%%%%%%%%%%%%%%%%%%%%%%%%%%%%%%%%%%
%%%%%% Last page in case there is no bibliography %%%%%%%%
% \label{Last Page}                           %%%%%%%%%%%%%
%%%%%%%%%%%%%%%%%%%%%%%%%%%%%%%%%%%%%%%%%%%%%%%%%%%%%%%%%%

\end{document}
